%% kapa: exam=1 round=20 year=2009 period=2 form=A subject=accounting subject_ko=회계학 seq=14 number=14
%% source: Q-Net 감정평가사 기출문제 (cst00309) artlSeq=5024178 file=extracted/1008/5024178/2009년 제20회 감정평가사 1차시험 기출문제/2교시.hwp
%% conversion: kordoc 4.24.51 -> tools/gen1.py (auto); review=proofread-sonnet; flags=none
\begin{question}{14}
(주)한국은 신기술 개발과 관련하여 20×8년과 20×9년에 각각 ￦120,000과 ￦60,000을 지출하였으며, 이 지출은 모두 무형자산 개발비의 인식조건을 충족한다. 당해 신기술은 20×9년 7월 1일부터 사용가능하게 되었다. 정액법으로 무형자산을 상각할 경우 20×9년 12월 31일 현재 대차대조표에 표시될 개발비의 장부금액은 얼마인가? (단, 개발비의 내용연수는 5년, 잔존가치는 없으며, 월할상각한다.)
\begin{choices}
  \item ￦144,000
  \item ￦149,750
  \item ￦150,000
  \item ￦161,750
  \item ￦162,000
\end{choices}
\end{question}
\answer{5}
